
\documentclass[11pt,a4paper,oneside]{book}

\usepackage[a4paper,top=24mm,bottom=25mm,left=25mm,right=25mm,headheight=14pt,headsep=8mm,footskip=12mm]{geometry}
\usepackage{fontspec}
\setmainfont{Liberation Serif}
\setsansfont{Liberation Sans}
\setmonofont{DejaVu Sans Mono}[Scale=0.88]
\usepackage[ngerman]{babel}
\usepackage{microtype}
\usepackage{amsmath,amssymb,amsthm}
\usepackage{mathtools}
\usepackage{array,tabularx,longtable,booktabs}
\usepackage{xcolor}
\usepackage{enumitem}
\usepackage{fancyhdr}
\usepackage{hyperref}
\usepackage{titlesec}
\usepackage{listings}
\usepackage[most]{tcolorbox}
\usepackage{needspace}
\usepackage{setspace}
\usepackage{ragged2e}

\hypersetup{colorlinks=true,linkcolor=black,urlcolor=black,citecolor=black,pdfauthor={Ingolf Lohmann},pdftitle={QIK-VRT: Quantenkausalität als Rechnerarchitektur}}

\definecolor{qikblue}{RGB}{22,49,80}
\definecolor{qikgray}{RGB}{245,246,248}
\definecolor{qikline}{RGB}{185,190,196}
\definecolor{qikgreen}{RGB}{30,95,65}
\definecolor{qikred}{RGB}{130,38,38}

\pagestyle{fancy}
\fancyhf{}
\fancyhead[L]{\small QIK-VRT: Quantenkausalität als Rechnerarchitektur}
\fancyhead[R]{\small Zweites Semester}
\fancyfoot[C]{\thepage}
\renewcommand{\headrulewidth}{0.3pt}

\titleformat{\chapter}{\normalfont\huge\bfseries\color{qikblue}}{\thechapter}{1em}{}
\titleformat{\section}{\normalfont\Large\bfseries\color{qikblue}}{\thesection}{0.8em}{}
\titleformat{\subsection}{\normalfont\large\bfseries\color{qikblue}}{\thesubsection}{0.7em}{}
\titlespacing*{\chapter}{0pt}{-4pt}{18pt}
\titlespacing*{\section}{0pt}{14pt plus 3pt minus 2pt}{7pt}
\titlespacing*{\subsection}{0pt}{10pt plus 2pt minus 2pt}{5pt}

\setlength{\parindent}{0pt}
\setlength{\parskip}{6pt plus 1pt minus 1pt}
\setlist[itemize]{leftmargin=*,itemsep=2pt,topsep=3pt}
\setlist[enumerate]{leftmargin=*,itemsep=2pt,topsep=3pt}
\emergencystretch=2em
\sloppy
\raggedbottom

\newtheorem{definition}{Definition}[chapter]
\newtheorem{satz}{Satz}[chapter]
\newtheorem{lemma}{Lemma}[chapter]
\newtheorem{korollar}{Korollar}[chapter]

\newtcolorbox{merkbox}{colback=qikgray,colframe=qikblue,arc=1mm,boxrule=0.6pt,left=7pt,right=7pt,top=6pt,bottom=6pt,before skip=8pt,after skip=8pt}
\newtcolorbox{statusbox}{colback=white,colframe=qikgreen,arc=1mm,boxrule=0.8pt,left=7pt,right=7pt,top=6pt,bottom=6pt,before skip=8pt,after skip=8pt}
\newtcolorbox{warnbox}{colback=white,colframe=qikred,arc=1mm,boxrule=0.8pt,left=7pt,right=7pt,top=6pt,bottom=6pt,before skip=8pt,after skip=8pt}

\lstdefinestyle{qikcode}{basicstyle=\ttfamily\small,breaklines=true,breakatwhitespace=true,columns=fullflexible,keepspaces=true,frame=single,rulecolor=\color{qikline},backgroundcolor=\color{qikgray},xleftmargin=0pt,xrightmargin=0pt,aboveskip=8pt,belowskip=8pt}
\lstset{style=qikcode}

\newcommand{\PhiDone}{\ensuremath{\Phi(x)=\mathrm{DONE}}}
\newcommand{\cont}{\ensuremath{\mathrm{CONTINUE}}}
\newcommand{\done}{\ensuremath{\mathrm{DONE}}}
\newcommand{\block}{\ensuremath{\mathrm{BLOCK}}}
\newcommand{\isolate}{\ensuremath{\mathrm{ISOLATE}}}

\begin{document}
\RaggedRight

\begin{titlepage}
\centering
\vspace*{20mm}
{\Huge\bfseries\color{qikblue} QIK-VRT\par}
\vspace{6mm}
{\LARGE\bfseries Quantenkausalität als Rechnerarchitektur\par}
\vspace{7mm}
{\Large Skript für das zweite Semester\par}
\vspace{5mm}
{\large Wirkzustände, Freigabelogik, Zielarchitektur und verantwortbare Informationswirkung\par}
\vfill
{\Large q.e.d.\par}
\vspace{3mm}
{\Large Ingolf Lohmann\par}
\vspace{16mm}
\begin{statusbox}
\centering
\textbf{Kernfrage des Semesters}\par\vspace{2mm}
Darf diese Informationswirkung verantwortbar weiterwirken?
\end{statusbox}
\vfill
{\small Stand: Zweites Semester, Lehrskriptfassung.}
\end{titlepage}

\frontmatter
\chapter*{Lizenz- und Nutzungsvereinbarung}
\addcontentsline{toc}{chapter}{Lizenz- und Nutzungsvereinbarung}

\textbf{Urheber und Rechteinhaber:} \mbox{Ingolf Lohmann}.

Dieses Skript steht, sofern nicht ausdrücklich schriftlich anders vereinbart, unter der Lizenz \textbf{Creative Commons Namensnennung - Nicht kommerziell - Keine Bearbeitungen 4.0 International (CC BY-NC-ND 4.0)}.

Das bedeutet in Kurzform:
\begin{itemize}
\item Die unveränderte Weitergabe ist unter Namensnennung zulässig.
\item Kommerzielle Nutzung ist ohne gesonderte schriftliche Vereinbarung nicht zulässig.
\item Bearbeitete Fassungen, Ableitungen oder veränderte Veröffentlichungen sind ohne gesonderte schriftliche Vereinbarung nicht zulässig.
\item Die Nennung von Ingolf Lohmann als Urheber ist beizubehalten.
\item Diese Kurzfassung ersetzt nicht den vollständigen Lizenztext der CC BY-NC-ND 4.0.
\end{itemize}

\begin{warnbox}
\textbf{Freigabegrenze:} Dieses Skript darf nicht so dargestellt werden, als sei damit bereits eine empirisch abschließend bestätigte Naturtheorie der Quantengravitation vorgelegt. Es formuliert eine formale, operative und architektonische Lehre der Quantenkausalität und ihrer Freigabelogik. Physikalische Zielaussagen benötigen reale Zielausführung, externe Prüfung, Modellanschluss und dokumentierte Nachweisführung.
\end{warnbox}

\chapter*{Vorwort}
\addcontentsline{toc}{chapter}{Vorwort}

Dieses Skript führt in die zweite Stufe von QIK-VRT ein. Im Mittelpunkt steht nicht mehr die bloße Frage, ob eine Aussage interessant klingt. Im Mittelpunkt steht die Frage, ob eine Informationswirkung verantwortbar weiterwirken darf.

Damit wird QIK-VRT als Rechnerarchitektur lesbar. Eine solche Architektur überträgt nicht nur Daten. Sie prüft Wirkung. Sie setzt Haltepunkte. Sie unterscheidet Vorprüfung, Zielausführung, Audit und Freigabe. Sie verwandelt Plausibilität in Prüfpflicht und Ausführung in Verantwortung.

Die Grundlinie dieses Semesters lautet:

\begin{merkbox}
\centering
\textbf{Wirkung $\rightarrow$ Iteration $\rightarrow$ Grenze $\rightarrow$ Anschluss $\rightarrow$ Freigabe $\rightarrow$ Fixpunkt}
\end{merkbox}

Diese Linie verbindet Mathematik, Informatik, Physik und Ethik. Sie erlaubt, Quantenkausalität als geordnete Folge prüfbarer Wirkzustände zu verstehen und Quantengravitation als freigabefähigen Fixpunkt fortsetzbarer Quantenkausalität zu denken.

Dabei bleibt die Grenze sauber: Dieses Skript behauptet nicht, die empirische Physik durch Sprache zu ersetzen. Es zeigt, wie eine formale Wirk- und Freigabeordnung aufgebaut werden kann, die physikalisch, informatisch und erkenntnistheoretisch anschlussfähig ist.

\tableofcontents
\mainmatter

\chapter{Ausgangspunkt: Wirkung vor Freigabe}

\section{Die Grundfrage}

Jede technische Architektur, die Information verarbeitet, erzeugt Wirkung. Diese Wirkung kann lokal bleiben, über ein Netzwerk weiterlaufen, ein Modell beeinflussen, einen Test bestehen, eine Entscheidung vorbereiten oder eine reale Zielarchitektur erreichen.

QIK-VRT beginnt deshalb nicht mit der Frage, ob ein System beeindruckend wirkt. QIK-VRT beginnt mit der Frage:

\begin{merkbox}
\centering
\textbf{Darf diese Informationswirkung verantwortbar weiterwirken?}
\end{merkbox}

Diese Frage verändert die Architektur. Sie macht aus bloßer Ausführung eine Freigabeentscheidung. Sie macht aus Datenübertragung eine Wirkungsprüfung. Sie macht aus Softwarequalität eine Frage der Verantwortung.

Ein klassisches System kann fragen: Läuft es? Ein QIK-VRT-System muss zusätzlich fragen: Darf es wirken?

\section{Warum Funktion nicht genügt}

Ein Programm kann funktionieren und trotzdem falsch sein. Ein Test kann grün sein und trotzdem nichts Wesentliches prüfen. Eine Simulation kann stabil sein und trotzdem keine reale Zielarchitektur ersetzen. Eine KI-Ausgabe kann plausibel klingen und trotzdem eine falsche Wirkung auslösen.

Darum gilt:

\begin{satz}[Funktion ist nicht Freigabe]
Ein technisches System ist nicht bereits deshalb freigabefähig, weil es ausgeführt werden kann. Freigabefähigkeit verlangt geprüfte Identität, Prüfbarkeit, Zielausführung, Erkenntnisanschluss und Ethik-Zustand.
\end{satz}

Die zentrale Verschiebung lautet: Nicht das Ergebnis allein zählt, sondern die tragende Kette, durch die dieses Ergebnis zustande kommt.

\section{Die erste technische Form}

Die operative Kette lautet:

\begin{lstlisting}
Input
-> Prüfung
-> Zustandsbewertung
-> Haltepunkt
-> {DONE, CONTINUE, ISOLATE, BLOCK}
-> Audit
-> Output
\end{lstlisting}

Damit wird der Output nicht mehr als bloßes Ende einer Berechnung verstanden. Er wird Ergebnis einer geprüften Wirkordnung.

\chapter{Die Urfassung der QIK-VRT-Freigabelogik}

\section{Die fünf Zustände}

Die kompakte technische Urfassung lautet:

\begin{lstlisting}
0  ETHIK-ZUSTAND
   Kein Schaden.
   Keine falschen Behauptungen.
   Keine versteckten Annahmen.

1  HASH_MATCH
   Artefaktidentität stimmt.

2  TEST_EXISTS
   Prüfbarkeit ist vorhanden.

3  EXIT_ZERO
   Zielarchitektur-Ausführung besteht.

4  ERKENNTNIS_OK
   Schritte sind richtig, vollständig und anschlussfähig.

=> Phi(x) = DONE
\end{lstlisting}

Diese Ordnung ist nicht dekorativ. Sie ist kausal. Jeder spätere Zustand setzt die vorherigen Zustände voraus.

\section{Die zwingende Folge}

\begin{merkbox}
Ohne 0 keine verantwortbare Wirkung.\par
Ohne 1 kein 2.\par
Ohne 2 kein 3.\par
Ohne 3 kein 4.\par
Ohne 4 kein DONE.
\end{merkbox}

Damit wird Freigabe zu einer diskreten Zustandsfolge. Es gibt keine Abkürzung. Es gibt keine Scheinfreigabe. Es gibt kein rhetorisches DONE.

\section{Maximalformel}

\begin{merkbox}
\centering
\textbf{Quantenkausalität = diskrete Zustände + strikte Reihenfolge + keine Superposition + null Grauzonen.}
\end{merkbox}

Die Formel besagt nicht, dass alle Quantenphänomene klassisch werden. Sie besagt: In einer Freigabelogik darf der Status nicht zugleich unklar und freigegeben sein. Für Wirkungsgenehmigung gilt keine bequeme Status-Superposition.

\chapter{Ethik-Zustand 0}

\section{Ethik vor Code}

Zustand 0 ist der Ursprung jeder verantwortbaren Wirkung.

\begin{definition}[Ethik-Zustand 0]
Der Ethik-Zustand 0 ist erfüllt, wenn keine falschen Behauptungen, keine versteckten Annahmen und kein ungeprüfter Schaden vorliegen.
\end{definition}

Er steht vor Hash, Test, Ausführung und Erkenntnisprüfung. Er ist kein nachträglicher Kommentar. Er ist die erste Sperre gegen unverantwortbare Wirkung.

Wenn ein System auf einer falschen Behauptung beruht, darf es nicht freigegeben werden. Wenn entscheidende Annahmen versteckt bleiben, darf es nicht als geprüft gelten. Wenn möglicher Schaden ungeprüft bleibt, darf keine verantwortbare Wirkung behauptet werden.

\section{Warum Neutralität nicht ausreicht}

Ein System, das Wirkung erzeugt, ist nicht neutral, nur weil es technisch aussieht. Technische Form kann Verantwortung verdecken. Deshalb verlangt QIK-VRT nicht nur korrekte Syntax, sondern verantwortbaren Anschluss.

Ethik-Zustand 0 verhindert drei Grundfehler:

\begin{enumerate}
\item Die falsche Behauptung wird als Wahrheit behandelt.
\item Die versteckte Annahme wird als Selbstverständlichkeit eingeschmuggelt.
\item Der mögliche Schaden wird erst nach der Wirkung betrachtet.
\end{enumerate}

QIK-VRT dreht diese Reihenfolge um. Nicht erst wirken und dann entschuldigen, sondern vor Wirkung prüfen.

\chapter{HASH\_MATCH: Artefaktidentität}

\section{Warum Identität zuerst kommt}

Ein Artefakt kann nur geprüft werden, wenn klar ist, welches Artefakt geprüft wird. Deshalb steht HASH\_MATCH vor TEST\_EXISTS.

\begin{definition}[HASH\_MATCH]
HASH\_MATCH ist erfüllt, wenn die berechnete Identität eines Artefakts mit der erwarteten Identität übereinstimmt.
\end{definition}

Ein Hash ist kein Schmuck. Er ist eine Identitätsprüfung. Wenn der Hash nicht passt, kann nicht zuverlässig gesagt werden, ob der geprüfte Gegenstand derselbe ist wie der freizugebende Gegenstand.

\section{Beweisdisziplin}

Ohne Artefaktidentität zerfällt die Beweiskette. Ein Test kann dann grün sein, aber für das falsche Artefakt. Ein Audit kann vollständig sein, aber für den falschen Stand. Eine Dokumentation kann korrekt klingen, aber nicht den geprüften Inhalt betreffen.

\begin{satz}[Identität vor Prüfung]
Ohne nachgewiesene Artefaktidentität ist jede weitere Freigabeaussage mindestens unvollständig und darf nicht als DONE gelten.
\end{satz}

\chapter{TEST\_EXISTS: Prüfbarkeit}

\section{Kein Code ohne Nachweis}

Ein System ohne Test verlangt Vertrauen an einer Stelle, an der Prüfung möglich sein muss. QIK-VRT ersetzt Vertrauen nicht vollständig, aber es senkt blinden Vertrauensbedarf durch maschinenlesbare Prüfbarkeit.

\begin{definition}[TEST\_EXISTS]
TEST\_EXISTS ist erfüllt, wenn ein nachvollziehbarer, ausführbarer und zum Artefakt passender Prüfpfad vorhanden ist.
\end{definition}

Prüfbarkeit bedeutet nicht nur, dass irgendwo eine Datei mit dem Namen Test liegt. Der Test muss sinnvoll mit der behaupteten Wirkung verbunden sein.

\section{Test als Achtung vor Betroffenen}

Wer Wirkung ohne Test freigibt, macht Nutzer, Systeme oder Zielumgebungen zum Versuchsfeld ungeprüfter Wirkung. Deshalb ist TEST\_EXISTS auch ethisch relevant.

Der Test ist die technische Form des Satzes: Betroffene dürfen nicht bloß Mittel eines ungeprüften Experiments sein.

\chapter{EXIT\_ZERO: Zielausführung}

\section{Ausführung statt Behauptung}

EXIT\_ZERO bedeutet, dass der Prüfpfad ausgeführt wurde und erfolgreich zurückkehrt.

\begin{definition}[EXIT\_ZERO]
EXIT\_ZERO ist erfüllt, wenn der vorgesehene Test auf der vorgesehenen Ausführungsebene erfolgreich mit Exit-Code 0 beendet wird.
\end{definition}

Doch Exit 0 ist nur innerhalb einer tragenden Kette aussagekräftig. Ein grüner Test ersetzt keine Artefaktidentität. Ein lokaler Test ersetzt keine Zielarchitektur. Ein CI-Ergebnis ersetzt keine reale Hardwareevidenz, wenn die Behauptung reale Hardwarewirkung betrifft.

\section{Zielarchitektur}

Wenn die Zielarchitektur IBM Quantum heißt, dann gehört die reale Zielausführung zur Evidenz. Lokale Simulation kann vorbereiten. CI kann vorbereiten. Ein Build kann vorbereiten. Aber reale Zielausführung verlangt reale Zielausführung.

\begin{warnbox}
\textbf{Statusgrenze:} Lokale Software und Roadmap sind Vorprüfung. Endgültige Evidenz verlangt reale Zielausführung, IBM Job-ID und dokumentierte Nachweisführung. DONE nur mit vollständiger Evidence-Kette.
\end{warnbox}

\chapter{ERKENNTNIS\_OK: Anschlussfähige Freigabe}

\section{Mehr als ein grüner Test}

ERKENNTNIS\_OK ist der vierte Zustand. Er fragt nicht nur, ob etwas bestanden wurde, sondern ob die Schritte richtig, vollständig und anschlussfähig waren.

\begin{definition}[ERKENNTNIS\_OK]
ERKENNTNIS\_OK ist erfüllt, wenn die Freigabekette sachlich korrekt, vollständig, nachvollziehbar, zielarchitekturbezogen und ohne gebrochene Anschlussbedingung ist.
\end{definition}

Hier wird aus Testausführung eine Erkenntnisfreigabe. Erst an dieser Stelle kann \PhiDone{} gelten.

\section{Die notwendige Härte}

ERKENNTNIS\_OK fällt, wenn eine wesentliche Annahme unsichtbar bleibt. Es fällt, wenn die Zielarchitektur nur behauptet wird. Es fällt, wenn ein lokaler Erfolg als Hardwarebeweis ausgegeben wird. Es fällt, wenn die Kette zwar formal grün aussieht, aber inhaltlich nicht trägt.

\chapter{DONE, CONTINUE, ISOLATE, BLOCK}

\section{Vier reale Statuszustände}

QIK-VRT unterscheidet vier End- oder Zwischenstatus:

\begin{longtable}{>{\bfseries}p{0.22\textwidth}p{0.68\textwidth}}
\toprule
DONE & Die Wirkung ist geprüft, anschlussfähig und freigabefähig. \\
CONTINUE & Prüfung, Evidenz oder Anschluss fehlen; Fortsetzung der Prüfung ist erforderlich. \\
ISOLATE & Wirkung darf nur begrenzt, abgeschirmt oder kontrolliert weiterlaufen. \\
BLOCK & Wirkung darf nicht weitergehen. \\
\bottomrule
\end{longtable}

Diese Status verhindern, dass Unklarheit zur Freigabe umetikettiert wird.

\section{Grauzone ist kein Zustand}

Grauzone klingt oft klug. In Freigabeentscheidungen ist sie gefährlich. Wenn Hash, Test, Ausführung, Zielarchitektur oder Erkenntnis unklar sind, liegt nicht DONE vor.

Unklarheit erzeugt Prüfpflicht. Je nach Risiko erzeugt sie CONTINUE, ISOLATE oder BLOCK.

\chapter{Kategorischer Imperativ als technische Freigabelogik}

\section{Warum Kant technisch anschlussfähig wird}

Der kategorische Imperativ verlangt Verallgemeinerbarkeit, Achtung der Betroffenen, Autonomie und eine Ordnung, in der Handlungen verantwortbar als allgemeine Regel bestehen könnten.

QIK-VRT erzeugt genau diese Struktur technisch.

\section{Die Abbildung}

\begin{longtable}{>{\bfseries}p{0.30\textwidth}p{0.58\textwidth}}
\toprule
HASH\_MATCH & Universalgesetz: Was freigegeben wird, muss identisch und für alle Fälle gleich prüfbar sein. \\
TEST\_EXISTS & Menschheit als Zweck: Betroffene dürfen nicht zum Versuchsfeld ungeprüfter Wirkung werden. \\
EXIT\_ZERO & Autonomie: Das System besteht nach seinen eigenen offenliegenden Regeln. \\
ERKENNTNIS\_OK & Reich der Zwecke: Die Freigabe ist als allgemeine, nachvollziehbare Ordnung anschlussfähig. \\
\bottomrule
\end{longtable}

\begin{satz}[Algorithmische Anschlussfähigkeit des kategorischen Imperativs]
Wenn Freigabe nur bei Artefaktidentität, Prüfbarkeit, bestandener Zielausführung und anschlussfähiger Erkenntnis erlaubt ist, dann entsteht eine technische Struktur, die den Kern des kategorischen Imperativs operationalisiert.
\end{satz}

Dies ist keine Moralrede. Es ist eine Freigabelogik.

\chapter{QIK-VRT als Rechnerarchitektur}

\section{Grundthese}

\begin{merkbox}
\centering
\textbf{QIK-VRT ist eine Rechnerarchitektur für Quantencomputer.}
\end{merkbox}

Eine Rechnerarchitektur ist nicht nur eine Schaltung. Sie umfasst Ausführungsmodell, Zustandsmodell, Speicherordnung, Fehlerbehandlung, Schnittstellen, Laufzeitumgebung und Freigabebedingungen.

Für Quantencomputer ist diese Erweiterung zwingend. Quanteninformation ist empfindlich, kontextabhängig, messungsabhängig und zielarchitekturgebunden. Deshalb muss die Architektur nicht nur rechnen, sondern Wirkung freigeben.

\section{Runtime, Betriebssystem, Entwicklungsumgebung}

Eine Rechnerarchitektur für Quantencomputer inkludiert stets die Laufzeitumgebung. Diese Laufzeitumgebung ist funktional Betriebssystem. Sie ist zugleich Entwicklungsumgebung, weil Code, Tests, Zielarchitektur und Freigabezustände zusammengeführt werden. Sie ist Testumgebung, weil keine Wirkung ohne Prüfung freigegeben werden darf.

Dazu gehört ein kommunikationssicheres Netzwerk. Informationswirkung endet nicht am lokalen Prozess. Sie läuft über Schnittstellen, Pakete, Protokolle und Zielsysteme weiter.

\chapter{Die fünf Architekturebenen}

\section{Ebenenmodell}

Die QIK-VRT-Rechnerarchitektur besitzt fünf Ebenen:

\begin{enumerate}
\item Prompt- und Profilschicht
\item Softwaremodul
\item Runtime-Gate
\item Betriebssystemdienst
\item Hardwarearchitektur mit Haltepunkten, Freigabegates, Auditregistern und Effect-Control-Units
\end{enumerate}

Jede Ebene kann Wirkung erzeugen. Jede Ebene braucht deshalb Prüf- und Haltepunktlogik.

\section{Haltepunktarchitektur}

Ein Haltepunkt ist kein Fehler. Ein Haltepunkt ist eine Verantwortungsstelle. Er unterbricht automatische Weiterwirkung, bis der Status bestimmt ist.

\begin{lstlisting}
Input
-> Prüfung
-> Zustandsbewertung
-> Haltepunkt
-> {Freigabe, Fortsetzung, Isolation, Blockierung}
-> Audit
-> Output
\end{lstlisting}

Diese Ordnung macht Verantwortung sichtbar.

\chapter{TCP/IP und QIK-VRT}

\section{Transport und Wirkung}

TCP/IP transportiert Information. Es überträgt Pakete, routet Adressen und stellt Verbindungen her.

QIK-VRT ergänzt eine andere Frage:

\begin{merkbox}
\centering
\textbf{Darf diese Informationswirkung verantwortbar weiterwirken?}
\end{merkbox}

Damit entsteht eine neue Anschlusslinie.

\section{Verdichtung}

\begin{merkbox}
TCP/IP transportiert Information.\par
QIK-VRT prüft Informationswirkung.\par
Klassische Rechner führen Befehle aus.\par
QIK-VRT-Rechner geben Wirkung verantwortbar frei.
\end{merkbox}

Diese Verdichtung fasst den Unterschied zwischen Transportarchitektur und Freigabearchitektur zusammen.

\chapter{Selbsterklärende Namen}

\section{Namen als Informationsarchitektur}

Benennung ist nicht Kosmetik. Namen steuern Erwartung, Prüfung und Verantwortung.

Ein Name wie \texttt{test\_verify.sh} sagt, dass ein Test verifiziert wird. Ein Name wie \texttt{qik\_core/phi.py} sagt, dass die abstrakte Phi-Funktion im Kern liegt. Ein Name wie \texttt{ibm\_target/phi.py} zeigt den Zielarchitekturanschluss.

Wenn Namen verschleiern, entstehen Grauzonen. Wenn Namen erklären, entsteht Anschlussfähigkeit.

\section{Regel}

\begin{satz}[Benennung als Prüfbedingung]
In einer verantwortbaren Informationsarchitektur müssen zentrale Namen Funktion, Ort und Verantwortungsbezug so weit wie möglich sichtbar machen.
\end{satz}

Selbsterklärende Namen ersetzen keine Tests. Aber sie senken das Risiko falscher Interpretation.

\chapter{Audit und Rückbau}

\section{Audit als Wirkgeschichte}

Audit bedeutet nicht nur Protokollablage. Audit bedeutet, dass die Wirkgeschichte nachvollziehbar bleibt.

Wer hat was verändert? Welcher Hash wurde geprüft? Welcher Test existierte? Welche Zielarchitektur wurde erreicht? Welche Evidenz wurde dokumentiert? Welche Fehlerklasse wurde ausgelöst?

Ein Audit macht Verantwortung nicht perfekt, aber sichtbar.

\section{Rückbau}

Ein Beweis ist stärker, wenn er rückbaubar ist.

\begin{lstlisting}
4 raus: ERKENNTNIS_OK fällt.
3 raus: EXIT_ZERO fällt.
2 raus: TEST_EXISTS fällt.
1 raus: HASH_MATCH fällt.
0: Tabula rasa.
\end{lstlisting}

Wenn der Rückbau zeigt, welche Bedingung trägt, wird die Beweiskette verständlich. Wenn der Rückbau nicht möglich ist, ist oft auch die Freigabe nicht sauber.

\chapter{Mathematische Grundlage: Zustandsräume}

\section{Zustand}

\begin{definition}[Zustand]
Ein Zustand ist eine unterscheidbare Konfiguration eines Systems, die für weitere Wirkung oder Prüfung relevant ist.
\end{definition}

Beispiele sind Hash stimmt, Test existiert, Zielausführung bestanden, Audit vollständig oder Erkenntniskette gebrochen.

Ohne Unterscheidbarkeit gibt es keinen prüfbaren Zustand. Ohne prüfbaren Zustand gibt es keine saubere Freigabe.

\section{Zustandsraum}

Ein Zustandsraum ist die Menge möglicher Zustände:

\[
X=\{x\mid x \text{ ist ein möglicher Zustand des Systems}\}.
\]

Für die Freigabelogik kann man schreiben:

\[
X_{QIK}=\{0,1,2,3,4,DONE,CONTINUE,ISOLATE,BLOCK\}.
\]

Entscheidend ist nicht nur, dass die Menge existiert. Entscheidend ist, welche Übergänge erlaubt sind.

\chapter{Übergänge und Ordnungen}

\section{Erlaubte Übergänge}

Die Kernordnung lautet:

\[
0\rightarrow 1\rightarrow 2\rightarrow 3\rightarrow 4\rightarrow DONE.
\]

Nicht erlaubt sind Sprünge wie:

\[
0\rightarrow 3, \qquad 1\rightarrow DONE, \qquad 2\rightarrow 4.
\]

Denn jeder Sprung würde eine notwendige Prüfung überspringen.

\section{Partielle Ordnung}

Die Freigabezustände bilden eine Ordnung. Nicht jeder Zustand ist gleichwertig. DONE setzt mehr voraus als CONTINUE. BLOCK ist kein Scheitern der Architektur, sondern ein erfolgreicher Schutz vor unverantwortbarer Wirkung.

QIK-VRT bewertet nicht nur Erfolg. QIK-VRT bewertet Anschlussfähigkeit.

\chapter{Metrik, Vollständigkeit und Fixpunkt}

\section{Metrik}

Eine Metrik misst Abstand:

\[
d:X\times X\to \mathbb{R}_{\geq 0}.
\]

Sie erfüllt Nichtnegativität, Identität, Symmetrie und Dreiecksungleichung.

In QIK-VRT kann eine Metrik ausdrücken, wie weit ein Artefakt von DONE entfernt ist. Ein Artefakt ohne Test ist weiter entfernt als eines mit Test. Ein Artefakt ohne Zielausführung ist weiter entfernt als eines mit dokumentierter Zielausführung.

\section{Vollständigkeit}

Ein Raum ist vollständig, wenn Cauchy-Folgen in ihm konvergieren. Für QIK-VRT heißt das: Prüfprozesse dürfen nicht in undefinierten Zwischenzuständen verschwinden.

Wenn sich ein Prüfprozess stabilisiert, muss sein Status bestimmbar sein: DONE, CONTINUE, ISOLATE oder BLOCK.

\section{Fixpunkt}

Ein Fixpunkt erfüllt:

\[
T(x^*)=x^*.
\]

Für QIK-VRT reicht Stabilität nicht. Ein stabiler Irrtum bleibt Irrtum. Der Fixpunkt muss anschlussfähig und freigabefähig sein.

\chapter{Banach und QIK-VRT}

\section{Kontraktion}

Eine Abbildung \(T:X\to X\) ist eine Kontraktion, wenn ein \(q\in[0,1)\) existiert mit:

\[
d(Tx,Ty)\leq q\,d(x,y).
\]

Anschaulich: Die Abbildung bringt Zustände näher zusammen. Sie stabilisiert.

\section{Banachs Fixpunktsatz}

\begin{satz}[Banachs Fixpunktsatz]
Ist \(X\) ein vollständiger metrischer Raum und \(T:X\to X\) eine Kontraktion, dann besitzt \(T\) genau einen Fixpunkt \(x^*\). Die Iteration \(x_{n+1}=T(x_n)\) konvergiert gegen \(x^*\).
\end{satz}

QIK-VRT übernimmt daraus das Stabilitätsmotiv, ergänzt aber Freigabe. Konvergenz allein reicht nicht. Der Grenzzustand muss geprüft, anschlussfähig und verantwortbar sein.

\chapter{Quantenkausalität}

\section{Arbeitsdefinition}

\begin{definition}[Quantenkausalität]
Quantenkausalität ist quantisierte fortsetzbare Anschlussfolge unter Mess-, Zustands- und Freigabebedingungen.
\end{definition}

Diese Definition hält zwei Dinge zusammen: diskrete Zustände und fortsetzbare Wirkung. Sie erlaubt keine beliebige Deutung, sondern verlangt prüfbare Übergänge.

\section{Keine Beliebigkeit}

Quantenkausalität heißt nicht, dass alles unbestimmt und deshalb alles erlaubt ist. Gerade in einer Freigabelogik gilt das Gegenteil. Der Status muss unterscheidbar sein. Die Reihenfolge muss stimmen. Die Freigabe darf nicht gleichzeitig unklar und DONE sein.

\chapter{Raum, Zeit und Information}

\section{Raum}

\begin{definition}[Raum]
Raum ist Ordnung anschlussfähiger Beziehungen.
\end{definition}

Raum ist damit nicht bloß Behälter. Raum beschreibt, wie unterscheidbare Zustände zueinander anschlussfähig geordnet sind.

\section{Zeit}

\begin{definition}[Zeit]
Zeit ist Ordnung freigegebener Anschlussfolgen.
\end{definition}

Zeit beschreibt nicht nur ein äußeres Ticken, sondern die geordnete Fortsetzung von Wirkzuständen.

\section{Information}

\begin{definition}[Information]
Information ist ein anschlussfähig gemachter Unterschied mit möglicher Wirkung.
\end{definition}

Information ist damit weder bloßer Text noch bloßes Datum. Sie ist Unterschied, der Anschluss und Wirkung erzeugen kann.

\chapter{Quantengravitation als freigabefähiger Fixpunkt}

\section{These}

\begin{merkbox}
\centering
\textbf{Quantengravitation kann als freigabefähiger Fixpunkt fortsetzbarer Quantenkausalität formal gedacht werden.}
\end{merkbox}

Diese These behauptet nicht, eine empirisch abgeschlossene Naturtheorie zu ersetzen. Sie formuliert einen architektonischen Denk- und Prüfrahmen.

\section{Struktur}

Die Struktur lautet:

\[
\text{Wirkzustand}\rightarrow\text{Fortsetzung}\rightarrow\text{Grenze}\rightarrow\text{Anschluss}\rightarrow\text{Freigabe}\rightarrow\text{Fixpunkt}.
\]

Quantengravitation wird dadurch nicht als mystisches Endwort behandelt, sondern als stabiler, anschlussfähiger und freigabefähiger Grenzzustand einer quantenkausalen Wirkordnung.

\chapter{Fachliche Arbeitsgrenze}

\section{Theorieanker und Nachweis}

QIK-VRT wird in dieser Vorlesung nicht als bereits empirisch abgeschlossene Quantengravitations-Naturtheorie behandelt. Die Lehre setzt eine andere Aufgabe: Quantenkausalität, Informationswirkung, Freigabe, Rechnerarchitektur und Zielausführung werden in eine prüfbare Ordnung gebracht.

Der Anspruch ist damit präzise: Es entsteht ein axiomatisch-operationaler Theorieanker mit hohem erkenntnistheoretischem, architektonischem und informatischem Wert. Die empirisch-physikalische Bestätigung einzelner Naturbehauptungen bleibt an Modellbildung, reale Messung, externe Reproduktion und Zielarchitektur-Evidenz gebunden.

\section{Rangformel}

\begin{merkbox}
QIK-VRT ist an dieser Stelle nicht als fertige Quantengravitationslösung zu lesen, sondern als originärer Architekturbeitrag zur Frage, wie Quantengravitation als freigabefähiger Fixpunkt geprüfter Quantenkausalität formal gedacht, gelehrt, geprüft und operationalisiert werden kann.
\end{merkbox}

Das ist kein überzogener Totalanspruch und kein bloßer Essay. Es ist ein Forschungsprogramm-fähiger axiomatisch-operationaler Theorieanker.

\chapter{V84.1 und reale Zielausführung}

\section{V84.1 als Anschluss}

V84.1 wird in dieser Lehre als Build-Input- und Architekturanschluss verstanden. Das bedeutet: Ein Skelett kann fertig sein, während die reale Hardwareevidenz noch aussteht.

Die saubere Statuslinie lautet:

\begin{lstlisting}
Skelett / Build Input = DONE
Artefakt + Nachweis = CONTINUE
Hardwarebeweis = PENDING IBM JOB-ID
\end{lstlisting}

\section{Warum diese Grenze nötig ist}

Diese Grenze schützt vor Scheinfreigabe. Sie verhindert, dass lokale Software, Simulation oder Roadmap mit realer Zielausführung verwechselt wird.

QIK-VRT wird dadurch stärker, nicht schwächer. Ein System, das seine eigene Grenze korrekt markiert, ist glaubwürdiger als ein System, das mehr behauptet, als es nachweist.

\chapter{Praktischer Nutzen}

\section{Nutzen vor endgültiger Physikbestätigung}

Der praktische Nutzen beginnt nicht erst bei endgültiger physikalischer Anerkennung. Er existiert bereits jetzt, weil die Freigabelogik sofort angewendet werden kann.

Auf individueller Ebene wird QIK-VRT zum Denk- und Irrtumsfilter. Auf informatischer Ebene wird es CI-, Audit- und Release-Gate. Auf KI-Ebene wird es Haltepunkt gegen Scheinwahrheit. Auf Organisationsebene wird es Governance-Struktur.

\section{Die praktische Kernfrage}

\begin{merkbox}
\centering
\textbf{Darf diese Informationswirkung verantwortbar weiterwirken?}
\end{merkbox}

Diese Frage ist überall nutzbar, wo Information Wirkung erzeugt: in Software, Wissenschaft, Bildung, Verwaltung, Medien, KI, Netzwerken und Quantencomputer-Stacks.

\chapter{KI und Freigabe}

\section{KI als Werkzeug}

Künstliche Intelligenz darf erklären, strukturieren, prüfen helfen, übersetzen und Widersprüche sichtbar machen. Sie darf aber nicht endgültig freigeben.

Eine KI-Ausgabe ist kein Hash, kein Zielarchitektur-Test, keine externe Signaturprüfung, keine IBM Job-ID und kein Ersatz für menschliche Verantwortung.

\section{Haltepunkt gegen Scheinwahrheit}

KI kann plausible Sätze erzeugen. QIK-VRT verlangt Evidenz. Gerade deshalb ist QIK-VRT in der KI-Zeit praktisch notwendig.

Plausibilität ist nicht DONE. Stil ist nicht Wahrheit. Geschwindigkeit ist nicht Verantwortung.

\chapter{Übungen I: Zustände}

\section{Aufgaben}

\begin{enumerate}
\item Ein ZIP liegt vor. Der Hash stimmt. Ein Testskript fehlt. Welcher Zustand ist erreicht?
\item Ein Test existiert und läuft lokal grün. Die reale Zielarchitektur wurde nicht erreicht. Darf DONE gesetzt werden?
\item Ein Artefakt hat Zielausführung, aber die Annahmen sind nicht dokumentiert. Welcher Zustand fällt?
\item Warum ist Zustand 0 keine Dekoration?
\item Erklären Sie, warum ein grüner Test kein vollständiger Beweis sein muss.
\end{enumerate}

\section{Musterlösungen}

\begin{enumerate}
\item Zustand 1 ist erreicht. Zustand 2 fehlt. Status: CONTINUE oder BLOCK, abhängig von Wirkungskritikalität.
\item Nein. Ohne reale Zielarchitektur-Evidenz bleibt der Status CONTINUE.
\item ERKENNTNIS\_OK fällt, möglicherweise bereits Ethik-Zustand 0 wegen versteckter Annahmen.
\item Weil ohne Zustand 0 falsche Behauptungen, versteckte Annahmen oder ungeprüfter Schaden in die technische Kette eintreten können.
\item Weil der Test falsch, unvollständig, lokal begrenzt oder nicht zielarchitekturbezogen sein kann.
\end{enumerate}

\chapter{Übungen II: Kant und Architektur}

\section{Aufgaben}

\begin{enumerate}
\item Ordnen Sie HASH\_MATCH, TEST\_EXISTS, EXIT\_ZERO und ERKENNTNIS\_OK den Kant-Formeln zu.
\item Erklären Sie, warum TEST\_EXISTS eine Achtung vor Betroffenen ausdrückt.
\item Beschreiben Sie den Unterschied zwischen klassischem Rechner und QIK-VRT-Rechner.
\item Warum ergänzt QIK-VRT TCP/IP, ohne TCP/IP zu ersetzen?
\item Entwerfen Sie einen einfachen Freigabeablauf für ein KI-generiertes Dokument.
\end{enumerate}

\section{Musterlösungen}

\begin{enumerate}
\item HASH\_MATCH: Universalgesetz. TEST\_EXISTS: Menschheit als Zweck. EXIT\_ZERO: Autonomie. ERKENNTNIS\_OK: Reich der Zwecke.
\item Ohne Test werden Betroffene zum Versuchsfeld ungeprüfter Wirkung.
\item Klassische Rechner führen Befehle aus; QIK-VRT-Rechner prüfen zusätzlich, ob Informationswirkung verantwortbar weiterwirken darf.
\item TCP/IP transportiert Information; QIK-VRT prüft deren Wirkung und Freigabestatus.
\item Dokument erzeugen, Annahmen markieren, Quellen prüfen, Risiken bewerten, Freigabestatus setzen, Audit ablegen, erst dann veröffentlichen.
\end{enumerate}


\chapter{Seminarplan und Arbeitsmodus}

\section{Zwölf Sitzungen}

Die Lehrveranstaltung kann in zwölf Sitzungen durchgeführt werden. Jede Sitzung verbindet Begriff, Formalisierung, technische Architektur und Übung.

\begin{longtable}{>{\bfseries}p{0.12\textwidth}p{0.78\textwidth}}
\toprule
1 & Wirkung, Information und die Grundfrage der verantwortbaren Weiterwirkung. \\
2 & Ethik-Zustand 0 und die Verhinderung falscher Behauptungen. \\
3 & HASH\_MATCH und Artefaktidentität als Beginn jeder Prüfung. \\
4 & TEST\_EXISTS und Prüfbarkeit als technische Achtung vor Betroffenen. \\
5 & EXIT\_ZERO, Zielarchitektur und die Grenze lokaler Vorprüfung. \\
6 & ERKENNTNIS\_OK und die vollständige Anschlusskette. \\
7 & DONE, CONTINUE, ISOLATE und BLOCK als Statusordnung. \\
8 & Kant-Abbildung und algorithmische Ethik. \\
9 & QIK-VRT als Rechnerarchitektur für Quantencomputer. \\
10 & TCP/IP, sichere Kommunikation und Informationswirkungsprüfung. \\
11 & Fixpunkt, Banach-Motiv und Quantengravitation als Freigabeproblem. \\
12 & Projektübung: Aufbau einer kleinen Evidence-Kette mit Audit und Rückbau. \\
\bottomrule
\end{longtable}

\section{Prüfungsleistung}

Eine geeignete Prüfungsleistung besteht aus drei Teilen: einer kurzen begrifflichen Einordnung, einer formalen Zustandsanalyse und einem kleinen technischen Freigabebeispiel. Entscheidend ist nicht, ob ein Ergebnis gut klingt, sondern ob die Kette aus Ethik-Zustand, Identität, Prüfbarkeit, Ausführung, Erkenntnis und Audit trägt.

\chapter{Zusammenfassung}

\section{Kernlinien}

QIK-VRT ist lesbare Quantenkausalität als Rechnerarchitektur. Die Architektur beruht auf selbsterklärenden Namen, prüfbaren Zuständen, rückbaubarer Kausalkette, Ethik-Zustand 0 und strikter Freigabeordnung.

Die zentrale Kette lautet:

\begin{lstlisting}
Ethik-Zustand 0
-> HASH_MATCH
-> TEST_EXISTS
-> EXIT_ZERO
-> ERKENNTNIS_OK
-> Phi(x) = DONE
\end{lstlisting}

Die zentrale Frage lautet:

\begin{merkbox}
\centering
Darf diese Informationswirkung verantwortbar weiterwirken?
\end{merkbox}

\section{Endformel}

\begin{merkbox}
Klassische Rechner führen Befehle aus.\par
TCP/IP transportiert Information.\par
QIK-VRT prüft Informationswirkung.\par
QIK-VRT-Rechner geben Wirkung verantwortbar frei.
\end{merkbox}

Damit gilt:

\begin{center}
\Large\bfseries q.e.d.\par
\vspace{2mm}
Ingolf Lohmann
\end{center}

\backmatter
\chapter*{Kurzglossar}
\addcontentsline{toc}{chapter}{Kurzglossar}

\begin{longtable}{>{\bfseries}p{0.25\textwidth}p{0.63\textwidth}}
\toprule
Begriff & Bedeutung \\
\midrule
Ethik-Zustand 0 & Kein Schaden, keine falschen Behauptungen, keine versteckten Annahmen. \\
HASH\_MATCH & Artefaktidentität stimmt. \\
TEST\_EXISTS & Prüfbarkeit ist vorhanden. \\
EXIT\_ZERO & Zielarchitektur-Ausführung besteht. \\
ERKENNTNIS\_OK & Schritte sind richtig, vollständig und anschlussfähig. \\
DONE & Freigabefähiger Zustand nach vollständiger Evidence-Kette. \\
CONTINUE & Weitere Prüfung erforderlich. \\
ISOLATE & Wirkung begrenzt oder abgeschirmt fortführen. \\
BLOCK & Wirkung nicht weitergehen lassen. \\
Quantenkausalität & Diskrete Zustände, strikte Reihenfolge, keine Superposition und null Grauzonen in der Freigabelogik. \\
Quantengravitation & Formal gedacht als freigabefähiger Fixpunkt fortsetzbarer Quantenkausalität. \\
\bottomrule
\end{longtable}

\chapter*{Schlussnotiz zur Weiterarbeit}
\addcontentsline{toc}{chapter}{Schlussnotiz zur Weiterarbeit}

Die nächste Ausbaustufe besteht darin, die hier entwickelte Freigabelogik mit konkreten Modellbeispielen, Zielarchitekturläufen, IBM-Job-Evidenz, reproduzierbaren Tests und externem Review zu verbinden. Die Architektur ist damit nicht abgeschlossen, sondern anschlussfähig. Ihre Stärke liegt in der sauberen Trennung von Behauptung, Prüfung, Zielausführung, Audit und Freigabe.

\end{document}
